% ===== PRÉAMBULE CANONIQUE — à coller VERBATIM en tête de CHAQUE .tex =====
\documentclass[11pt,a4paper]{article}
\usepackage[utf8]{inputenc}\usepackage[T1]{fontenc}\usepackage{lmodern}
\usepackage[french]{babel}
\usepackage{amsmath,amssymb,mathtools}\usepackage{icomma}
\usepackage{siunitx}\sisetup{locale=FR}  % \qty{}{}, \si{}, \ang{}
\usepackage{xcolor}\usepackage{colortbl}
\usepackage{tikz}\usepackage{pgfplots}\pgfplotsset{compat=1.18}
\usepackage[siunitx,europeanresistors]{circuitikz} % circuits (élec.)
\usepackage[most]{tcolorbox}\usepackage{enumitem}\usepackage{array,booktabs}
\usepackage[a4paper,margin=2.2cm]{geometry}
\usetikzlibrary{arrows.meta,calc,angles,quotes,positioning,patterns,%
  backgrounds,shapes.geometric,decorations.pathmorphing,decorations.markings,intersections}
\definecolor{primary}{RGB}{222,49,99}\definecolor{primarydark}{RGB}{150,22,55}
\definecolor{defcol}{RGB}{0,114,178}\definecolor{methcol}{RGB}{52,92,148}
\definecolor{excol}{RGB}{0,135,90}\definecolor{attncol}{RGB}{200,40,55}
\tcbset{boxbase/.style={enhanced,breakable,boxrule=0.9pt,arc=2.5pt,
  left=3mm,right=3mm,top=2.3mm,bottom=2.3mm,fonttitle=\bfseries,coltitle=white}}
% --- boîtes TUTORIEL (noms EXACTS, ne pas renommer) ---
\newtcolorbox{cle}[1][]{boxbase,colback=defcol!6,colframe=defcol,
  title={\textsc{À retenir}\ifx\relax#1\relax\relax\else\textmd{ : \itshape #1}\fi}}
\newtcolorbox{methode}[1][]{boxbase,colback=methcol!7,colframe=methcol,sharp corners,
  title={\textsc{Méthode}\ifx\relax#1\relax\relax\else\textmd{ --- \itshape #1}\fi}}
\newtcolorbox{exemplebox}[1][]{boxbase,colback=excol!5,colframe=excol,
  title={\textsc{Exemple}\ifx\relax#1\relax\relax\else\textmd{ : \itshape #1}\fi}}
\newtcolorbox{piege}[1][]{boxbase,colback=attncol!6,colframe=attncol,
  title={\textsc{Piège}\ifx\relax#1\relax\relax\else\textmd{ : \itshape #1}\fi}}
\newtcolorbox{remarque}[1][]{boxbase,colback=black!4,colframe=black!45,
  title={\textsc{Remarque}\ifx\relax#1\relax\relax\else\textmd{ : \itshape #1}\fi}}
% --- boîtes EXERCICE / CORRIGÉ (noms EXACTS, auto-comptées) ---
\newtcolorbox[auto counter]{exo}[1][]{boxbase,colback=primary!4,colframe=primary,
  title={\textsc{Exercice}~\thetcbcounter\ifx\relax#1\relax\relax\else\textmd{ --- \itshape #1}\fi}}
\newtcolorbox[auto counter]{corrige}[1][]{boxbase,colback=excol!4,colframe=excol,
  title={\textsc{Corrigé}~\thetcbcounter\ifx\relax#1\relax\relax\else\textmd{ --- \itshape #1}\fi}}
\newcommand{\vv}[1]{\vec{#1}}\newcommand{\ds}{\displaystyle}
\newcommand{\R}{\mathbb{R}}\newcommand{\N}{\mathbb{N}}\newcommand{\Z}{\mathbb{Z}}
% =========================================================================
\begin{document}

\begin{exo}[une réflexion totale qui trie les couleurs]
Un faisceau de lumière blanche se propage dans un verre dispersif ($n_{\text{rouge}}=1{,}51$,
$n_{\text{bleu}}=1{,}53$) et frappe la surface verre/air sous un angle d'incidence
$\hat{\imath}=\ang{41.0}$.
\begin{enumerate}[label=\alph*)]
  \item Calcule l'angle limite verre/air \emph{pour le rouge} puis \emph{pour le bleu}.
  \item Pour chaque couleur, dis si le rayon \textbf{sort} dans l'air (réfraction) ou subit une
        \textbf{réflexion totale}. \emph{(Indication : compare $n\sin\hat{\imath}$ à $1$.)}
  \item Quelle est la couleur dominante du faisceau \emph{transmis} dans l'air ? Et celle du faisceau
        \emph{réfléchi} dans le verre ?
\end{enumerate}
\end{exo}

\begin{exo}[d'où vient la formule de la profondeur apparente ?]
On veut établir $h'=\dfrac{n_2}{n_1}h$ pour un objet $O$ à la profondeur $h$ dans l'eau ($n_1$), vu
presque à la verticale dans l'air ($n_2$). On considère un rayon issu de $O$ faisant un \emph{petit}
angle $\hat{\imath}_1$ avec la verticale ; il atteint la surface à une distance horizontale $x$, se
réfracte sous un petit angle $\hat{\imath}_2$, et son prolongement semble venir de $O'$ à la
profondeur $h'$.
\begin{enumerate}[label=\alph*)]
  \item Pour de petits angles, $\tan\theta\approx\sin\theta$. Écris $x$ en fonction de $h,\hat{\imath}_1$
        puis de $h',\hat{\imath}_2$, et déduis-en $\dfrac{h'}{h}=\dfrac{\hat{\imath}_1}{\hat{\imath}_2}$.
  \item À l'aide de Snell–Descartes (petits angles : $n_1\hat{\imath}_1\approx n_2\hat{\imath}_2$),
        conclus que $h'=\dfrac{n_2}{n_1}h$.
  \item Application : un poisson est à $h=\qty{1.5}{m}$ dans l'eau ($n_1=1{,}33$). À quelle profondeur
        le voit-on ? De combien un pêcheur doit-il viser \emph{plus bas} ?
\end{enumerate}
\end{exo}

\begin{exo}[dispersion par un prisme]
Un rayon de lumière blanche entre \textbf{perpendiculairement} à la première face d'un prisme d'angle
au sommet $A=\ang{30}$ ($n_{\text{rouge}}=1{,}51$, $n_{\text{bleu}}=1{,}53$, prisme dans l'air). En
entrant sans être dévié, il frappe la seconde face sous un angle d'incidence interne égal à $A$.
\begin{center}
\begin{tikzpicture}[scale=1.0,>=Stealth]
  \coordinate (S) at (0,2.3); \coordinate (L) at (-0.9,0); \coordinate (R) at (2.6,0);
  \fill[defcol!18] (S)--(L)--(R)--cycle;
  \draw[black!70,line width=0.9pt] (S)--(L)--(R)--cycle;
  \node[above,font=\scriptsize] at (S){$A=\ang{30}$};
  \node[font=\scriptsize] at (0.55,0.55){prisme};
  \draw[->,line width=1pt] (-2.4,1.1)--(-0.42,1.1);
  \node[above,font=\scriptsize] at (-1.6,1.12){blanc};
  \coordinate (X) at (0.92,1.1);
  \draw[line width=0.8pt,black!35] (-0.42,1.1)--(X);
  \draw[->,line width=1pt,red] (X)--(3.0,0.55);
  \draw[->,line width=1pt,blue] (X)--(3.0,0.15);
  \fill (X) circle (0.8pt);
\end{tikzpicture}
\end{center}
\begin{enumerate}[label=\alph*)]
  \item Calcule l'angle sous lequel \emph{sort} le rouge, puis le bleu, à la seconde face.
  \item La déviation à cette face vaut $D=(\text{angle de sortie})-A$. Calcule $D_{\text{rouge}}$ et
        $D_{\text{bleu}}$, puis l'écart angulaire $\Delta D$ entre les deux couleurs.
  \item Quelle couleur est la plus déviée par le prisme ?
\end{enumerate}
\end{exo}

\end{document}
